\section{Conclusion}\label{sec:conclusion}

% Role: restate the question, the answer, and what to watch for as AI diffuses.

This paper asks whether AI-induced firm formation can offset the employment that large incumbents shed as they adopt AI.
We model the AI shock as a combination of displacement at incumbents and lower fixed costs of entry and operation, discipline the second channel using the cloud-computing era, and quantify the conditions for offset through calibration and Monte Carlo simulation.

\paragraph{Main findings.}
Across 2,000 Monte Carlo draws spanning plausible parameter ranges, the probability that AI raises net employment is 81.7\%.
The result is conditional: offset depends on the labor-AI substitution elasticity $\sigma_{\text{ai}}$, which explains 71\% of the variance in employment outcomes, and much of it reflects the absence of displacement rather than compensation for it.
Entry-cost reduction---the mechanism most analogous to cloud computing---explains only 3\%.

Below the displacement threshold ($\sigma_{\text{ai}} < 2.78$), AI is labor-complementary and offset is nearly universal (99.8\%).
Above the threshold, the offset probability falls to 74.9\% and depends heavily on whether AI-native entrants are more productive than their counterfactual selves.
When $\sigma_{\text{ai}} > 3.5$ and the entrant productivity advantage is weak, the offset probability collapses to 27.8\%.

Early empirical evidence is cautiously optimistic: the Census Bureau's BTOS (October 2026) shows that only 2\% of AI-adopting establishments reduced employment due to AI \citep{btos2026oct}, suggesting $\sigma_{\text{ai}}$ is currently below the threshold.
However, leading AI application firms are markedly employment-light---Cursor generates more than twenty times the revenue per employee of cloud-era SaaS firms---and business applications with planned wages are falling even as total applications reach record highs.

\paragraph{What to watch.}
The single most informative empirical signature is employment change at AI-adopting establishments, which directly reveals whether $\sigma_{\text{ai}}$ is above or below the displacement threshold.
Establishment-level data from the Business Trends and Outlook Survey (United States) and the Community Survey on ICT Usage (Europe) can track this as AI diffuses.
A shift from the current 2\% reporting employment reductions toward higher rates would signal that the labor-complementary regime is ending.

\paragraph{Broader implications.}
The standard approach to measuring AI's labor-market effects focuses on task exposure within a fixed set of firms.
Our analysis suggests this approach understates uncertainty: the net effect of AI on employment depends as much on firm creation as on within-firm automation.
Monitoring and forecasting should therefore track business formation, entrant productivity, and entrant employment alongside task automation rates.
The cloud era---which saw concentration rise despite cheaper computing---is a cautionary precedent: even favorable $\sigma_{\text{ai}}$ does not prevent concentration if structural headwinds suppress entry, and if $\sigma_{\text{ai}}$ turns out to be above the threshold, the same headwinds would turn offset into net job loss.
